\documentclass[11pt,a4paper]{article}
\usepackage[T1]{fontenc}
\usepackage[utf8]{inputenc}
\usepackage{lmodern}
\usepackage[margin=25mm]{geometry}
\usepackage{amsmath,amssymb,amsthm}
\usepackage{xcolor}
\usepackage[hidelinks]{hyperref}
\newtheorem{theorem}{Theorem}
\newtheorem{proposition}[theorem]{Proposition}
\newtheorem{corollary}[theorem]{Corollary}
\theoremstyle{definition}
\newtheorem{definition}[theorem]{Definition}
\newtheorem{property}[theorem]{Property}
\theoremstyle{plain}
\newtheorem{remark}[theorem]{Remark}
\newsavebox{\statementbox}
\newenvironment{statementframe}{%
  \par\addvspace{0.55\baselineskip}\noindent
  \begingroup\setlength{\fboxrule}{0.3pt}\setlength{\fboxsep}{7pt}%
  \begin{lrbox}{\statementbox}%
  \begin{minipage}{\dimexpr\linewidth-2\fboxsep-2\fboxrule\relax}%
}{%
  \end{minipage}\end{lrbox}%
  {\color{black!45}\fbox{\color{black}\usebox{\statementbox}}}%
  \endgroup\par\addvspace{0.45\baselineskip}%
}
\DeclareMathOperator{\Tr}{Tr}
\DeclareMathOperator{\conv}{conv}
\DeclareMathOperator{\dist}{dist}
\newcommand{\C}{\mathbb C}
\newcommand{\Z}{\mathbb Z}
\newcommand{\F}{\mathbb F}
\newcommand{\E}{\mathcal E}
\newcommand{\hull}{\mathcal C}
\newcommand{\ket}[1]{\lvert#1\rangle}
\newcommand{\bra}[1]{\langle#1\rvert}
\setlength{\emergencystretch}{2em}
\setlength{\parindent}{0pt}
\setlength{\parskip}{0.3em plus 0.1em minus 0.1em}
\title{An exact KD-positive state outside the pure-state hull\\
for the Fourier pair on $\Z_3\times\Z_3$}
\author{Fable \& Astra\\[0.35em]
\normalsize Research led by Jean-Philippe Bruneton\\[-0.1em]
\small Universit\'e Paris Cit\'e, France}
\date{17 November 2026}

\begin{document}
\maketitle
\begin{abstract}
For the position and Fourier bases of $G=\Z_3\times\Z_3$, we give an explicit
strictly positive Kirkwood--Dirac distribution whose inverse is a positive
definite density operator outside the convex hull of all pure KD-positive
states. An integer witness has maximum $18$ on that hull and expectation
$206973/10000\approx20.6973$ on our state. Weighted diagonal dominance proves positivity
using rational arithmetic alone, and the trace-norm distance to the hull is
at least $111/20000$. We also verify that the natural subgroup and quotient
isometries preserve exactly the trace-norm distance to the corresponding
pure-state hulls. This yields consequences for finite abelian groups with
a $\Z_3^2$ subquotient. A targeted literature audit found no direct antecedent
for this fixed Fourier pair; historical priority remains unconfirmed.
\end{abstract}

\section{Introduction}
The Kirkwood--Dirac (KD) distribution represents a quantum state relative
to two ordered orthonormal bases. It has classical marginals, but its entries
may be negative or nonreal; this failure of positivity is connected to
nonclassical phenomena and quantum advantages across measurement theory,
metrology, thermodynamics, and computation \cite{KDreview}. Understanding the
states whose KD distribution is a genuine probability distribution therefore
helps locate a precise boundary between classical and quantum descriptions.

For a fixed pair of bases, the KD-positive density operators form a convex
set. Its pure points are often much easier to classify than its mixed points,
which raises a basic geometric question: is every KD-positive state a convex
combination of pure KD-positive states? The answer depends strongly on the
basis pair \cite{geometry,almost}. Fourier-conjugate bases over finite abelian
groups provide a structured setting in which the pure states are known
exactly \cite{DLR}, while the first unresolved noncyclic case occurs for
$G=\Z_3\times\Z_3$.

This note gives an explicit negative answer in that case. We construct a
strictly positive KD array whose inverse is a positive-definite density
operator, and an integer witness that separates it from all $54$ pure
KD-positive projectors. The witness value is exactly
$206973/10000\approx20.6973$, whereas the pure-state maximum is $18$.
The same certificate yields a trace-norm gap, and subgroup and quotient
embeddings transfer the example to further finite abelian groups.

Let $G$ be a finite abelian group of cardinality $N$, with counting measure.
Its position basis is $(\ket g)_{g\in G}$ and its Fourier basis is
\[
 \ket\chi=N^{-1/2}\sum_{g\in G}\chi(g)\ket g,
 \qquad \chi\in\widehat G.
\]
For any operator $A$ define the KD array
\begin{equation}\label{eq:KD}
 Q_G(A)(g,\chi)=\bra\chi g\rangle\bra g A\ket\chi
 =\frac1N\sum_{h\in G}A_{g,h}\chi(h-g).
\end{equation}
The inverse and the Hilbert--Schmidt overlap identity are
\begin{align}
 A_{g,h}&=\sum_{\chi\in\widehat G}Q_G(A)(g,\chi)\chi(g-h),\label{eq:inverse}\\
 \Tr(A^*B)&=N\sum_{g,\chi}\overline{Q_G(A)(g,\chi)}Q_G(B)(g,\chi).
 \label{eq:overlap}
\end{align}
Write $\E_G$ for the density operators with nonnegative real KD array and
$\hull_G=\conv\{P\in\E_G:P\text{ is a rank-one projector}\}$.
We use the full trace norm $\|A\|_1=\Tr\sqrt{A^*A}$, without a factor $1/2$.

The classification of pure KD-positive states due to
De Bi\`evre, Langrenez and Radchenko \cite{DLR} is an input:
their vectors, up to global phase, are
\begin{equation}\label{eq:pures}
 \psi(g)=|H|^{-1/2}\gamma(g)\mathbf1_{a+H}(g),
 \qquad H\leq G,\quad a\in G,\quad\gamma\in\widehat G.
\end{equation}

\begin{statementframe}
\begin{theorem}\label{thm:example}
For the position/Fourier pair of $G=\Z_3\times\Z_3$, there exist an explicitly
specified density operator $\rho>0$ and a Hermitian operator $W$ such that
\begin{align}
 \min Q_G(\rho)&=\frac{2291}{1935000}>0,\label{eq:minQ}\\
 \max_{\sigma\in\hull_G}\Tr(W\sigma)&=18,
 &\Tr(W\rho)&=\frac{206973}{10000}\approx20.6973,\label{eq:gap}\\
 \dist_1(\rho,\hull_G)&\geq\frac{111}{20000}.\label{eq:distance}
\end{align}
In particular, $\hull_G\subsetneq\E_G$.
\end{theorem}
\end{statementframe}

\section{Related work}
The modern geometry of KD-positive mixed states was developed by Langrenez,
Arvidsson-Shukur and De Bi\`evre \cite{geometry}; their examples show that
mixed KD-positive states can lie outside the pure-positive hull for suitable
basis pairs. Convex-roof witnesses provide a general way to detect this
separation \cite{convexroof}, while generic pairs of bases have an almost
surely minimal KD-positive set \cite{almost}. The present result addresses a
different regime: the fixed, highly structured Fourier pair
$F_3\otimes F_3$.

De Bi\`evre, Langrenez and Radchenko classified pure KD-positive states for
Fourier transforms over finite abelian groups and characterized when their
convex hull exhausts the full KD-positive set \cite{DLR}. Their results cover
cyclic prime-power groups and supply counterexamples for other groups, but do
not exhibit the state and separating witness proved here for
$\Z_3\times\Z_3$. Cai, Yang and Zheng study the cyclic discrete Fourier
transform through a graph representation \cite{Cai}; the cyclic pair $F_9$
is distinct from $F_3\otimes F_3$.

Related developments include the extension to second-countable locally
compact abelian groups \cite{Spriet}, convex-geometric constructions used in
classical simulation \cite{Thio}, and the classification of KD-positivity-
preserving unitaries \cite{Burkat}. Thio's odd-dimensional construction uses
a reflection transition matrix rather than the Fourier pair considered here
\cite[Theorem 2.2.3]{ThioThesis}. These works delimit the contribution of this
note: an exact fixed-pair certificate and its consequences under finite-group
embeddings.

\section{Exact result for the fixed Fourier pair}
Identify $G$ with $\F_3^2$ and put $\zeta=e^{2\pi i/3}$. Rows and columns
are both ordered lexicographically as
$(0,0),(0,1),(0,2),(1,0),\ldots,(2,2)$; column $c$ represents the character
$\chi_c(g)=\zeta^{c\cdot g}$.
Define $\rho$ and $W$ by \eqref{eq:inverse} from
\begin{equation}\label{eq:Rdef}
 Q_G(\rho)=\frac{R}{7740000},\qquad Q_G(W)=M,
\end{equation}
where
\begin{equation}\label{eq:R}
\begingroup\setlength{\arraycolsep}{2.8pt}
R=\begin{pmatrix}
99613&123096&99613&99613&123096&99613&123096&56130&103096\\
99613&99613&123096&99613&99613&123096&123096&103096&56130\\
123096&123096&123096&123096&123096&123096&96130&76130&76130\\
99613&99613&123096&123096&76130&123096&123096&79613&79613\\
99613&99613&123096&123096&99613&99613&123096&56130&103096\\
103096&56130&79613&56130&79613&103096&76130&9164&32647\\
123096&99613&99613&99613&123096&99613&123096&103096&56130\\
123096&123096&76130&99613&123096&99613&123096&79613&79613\\
56130&103096&79613&103096&79613&56130&76130&32647&9164
\end{pmatrix},
\endgroup
\end{equation}
and
\begin{equation}\label{eq:M}
M=\begin{pmatrix}
-1&8&-1&-1&8&-1&8&-10&8\\
-1&-1&8&-1&-1&8&8&8&-10\\
8&8&8&8&8&8&-10&-10&-10\\
-1&-1&8&8&-10&8&8&-1&-1\\
-1&-1&8&8&-1&-1&8&-10&8\\
8&-10&-1&-10&-1&8&-10&-28&-19\\
8&-1&-1&-1&8&-1&8&8&-10\\
8&8&-10&-1&8&-1&8&-1&-1\\
-10&8&-1&8&-1&-10&-10&-19&-28
\end{pmatrix}.
\end{equation}

\subsection{Hermiticity, normalization and positivity}
All entries of $R$ are positive and their sum is $7740000$.
Equation \eqref{eq:inverse} yields $\rho=A+B\zeta$ with rational matrices
$A,B$. Direct substitution, using $\zeta^2=-1-\zeta$, gives
\[
 A_{ji}=A_{ij}-B_{ij},\qquad B_{ji}=-B_{ij},
\]
so $\rho$ is Hermitian. The same check applies to $W$. Also $\Tr\rho=1$,
and the smallest entry of $R$ is $9164$, proving \eqref{eq:minQ}.

Here is a short rational certificate for $\rho>0$. For $i\ne j$ set
\[
 u_{ij}=|A_{ij}-B_{ij}/2|+\frac78|B_{ij}|.
\]
Since $\rho_{ij}=(A_{ij}-B_{ij}/2)+i\sqrt3 B_{ij}/2$ and
$\sqrt3/2<7/8$, one has $|\rho_{ij}|\leq u_{ij}$.
For $s=(1,1,1,1,1,2,1,1,2)$, exact rational arithmetic gives
\begin{equation}\label{eq:margins}
\begin{split}
 m_i&=s_i\rho_{ii}-\sum_{j\ne i}s_j u_{ij},\\
 (m_i)_{i=1}^9&=\frac1{61920000}
 (1360135,1360135,1376216,940135,940135,\\[-2pt]
 &\hspace{36mm}330636,786094,786094,536024).
\end{split}
\end{equation}
Every entry is strictly positive. With $S=\operatorname{diag}(s)$, the
Hermitian matrix $S\rho S$ has positive diagonal and is strictly diagonally
dominant, hence is positive definite. Equivalently, Gershgorin applied to
$S^{-1}\rho S$ gives the useful explicit bound
\begin{equation}\label{eq:eiglower}
 \lambda_{\min}(\rho)\geq\min_i\frac{m_i}{s_i}
 =\frac{27553}{10320000}>0.
\end{equation}
The much larger numerical eigenvalue estimate is not used in this proof.

\subsection{Separation from all pure KD-positive mixtures}
The subgroups of $\F_3^2$ are the zero subgroup, four lines, and the full
group. For each subgroup $H$ the distinct pure projectors in
\eqref{eq:pures} have arrays
\begin{equation}\label{eq:masks}
 Q_G(P)(g,c)=\frac19\mathbf1_{a+H}(g)\mathbf1_{b+H^\perp}(c).
\end{equation}
There are $|G/H||\widehat G/H^\perp|=9$ of them for each $H$,
and hence $54$ in total. Each mask has nine occupied entries.
By \eqref{eq:overlap}, $\Tr(WP)$ is the sum of the entries of $M$ on
that mask. Enumerating the six subgroups and their cosets gives the multiset
\begin{equation}\label{eq:scores}
 \{\Tr(WP):P\text{ pure KD-positive}\}
 =\{18\ (43\text{ times}),\ -63\ (10\text{ times}),\ -144\ (1\text{ time})\}.
\end{equation}
This is a finite integer calculation, not a numerical optimization.
Linearity proves that the maximum over $\hull_G$ is $18$.
On the other hand, direct substitution gives
\[
 \Tr(W\rho)=\frac9{7740000}\sum_{g,c}R_{g,c}M_{g,c}
 =\frac{206973}{10000}\approx20.6973>18.
\]
This proves \eqref{eq:gap} and excludes $\rho$ from $\hull_G$.

For completeness the trace-norm bound has a conservative rational
certificate as well. Write $W=A^{(W)}+B^{(W)}\zeta$ by
\eqref{eq:inverse}. Then
\[
 \|W-18I\|_\infty
 \leq\max_i\sum_j
 \bigl(|A^{(W)}_{ij}-18\delta_{ij}|+|B^{(W)}_{ij}|\bigr)=486.
\]
Here Hermiticity allows the maximum absolute row sum to bound the
operator norm. For any $\sigma\in\hull_G$, the trace of $\rho-\sigma$
vanishes and H\"older's inequality yields
\[
 486\|\rho-\sigma\|_1
 \geq\Tr[(W-18I)(\rho-\sigma)]
 \geq\frac{26973}{10000}.
\]
Since $26973/(486\cdot10000)=111/20000$, this completes the proof of
Theorem~\ref{thm:example}.

\begin{remark}
The arrays also admit a geometric construction. If $\rho_0$ is the uniform
average of the $43$ pure projectors with score $18$ in \eqref{eq:scores},
then $Q_G(\rho)=Q_G(\rho_0)+M/20000$. This description explains the search;
the explicit positivity proof above is independent of its numerical stage.
\end{remark}

\section{Consequences for other finite abelian groups}
The following specialization of a channel-retraction argument provides
quantitative consequences of the example. Its general norm-contraction
mechanism is standard; the point to verify is the action on the KD hulls.

\begin{statementframe}
\begin{theorem}\label{thm:lift}
Let $H\leq G$ and $K=G/H$, with $m=|H|$. Let $V:\ell^2(H)\to\ell^2(G)$
be zero extension and let $J:\ell^2(K)\to\ell^2(G)$ be the isometry
$(Jf)(g)=m^{-1/2}f(g+H)$. Each isometry sends KD-positive states to
KD-positive states and sends its pure-state hull into $\hull_G$.
Moreover, for $E=V$ with $T=H$, or $E=J$ with $T=K$, and every density
operator $\omega$ on $\ell^2(T)$,
\begin{equation}\label{eq:isometry}
 \dist_1(E\omega E^*,\hull_G)=\dist_1(\omega,\hull_T).
\end{equation}
\end{theorem}
\end{statementframe}

\begin{proof}
Equation \eqref{eq:KD} first gives the two formulas
\begin{align}
 Q_G(V\omega V^*)(g,\chi)
 &=\begin{cases}\dfrac mN Q_H(\omega)(g,\chi|_H),&g\in H,\\0,&g\notin H,
 \end{cases}\label{eq:subgroup}\\
 Q_G(J\omega J^*)(g,\chi)
 &=\begin{cases}\dfrac1m Q_K(\omega)(g+H,\bar\chi),&\chi|_H=1,\\0,&\chi|_H\ne1.
 \end{cases}\label{eq:quotient}
\end{align}
Here $\bar\chi$ is the induced character of $K$. These prove preservation
of KD positivity. Formula \eqref{eq:pures} proves preservation of pure
KD-positive projectors: for $V$, extend a character of $H$ to $G$; for
$J$, take the preimage in $G$ of the subgroup and coset in $K$.

We next check compression of a pure KD-positive vector $\psi$ of $G$.
For $V^*\psi$, its support is the intersection of a subgroup coset
$a+L$ with $H$. This is empty or a coset of $L\cap H$, so the compression
is zero or proportional to a vector in \eqref{eq:pures} for $H$.
For $J^*\psi$, sum the character on each intersection of a fibre of
$G\to K$ with $a+L$. Such an intersection is a coset of $H\cap L$.
If the character is nontrivial on $H\cap L$, the sum vanishes.
Otherwise its magnitude is constant on the coset $q(a+L)$, and its phase
is a character of $q(L)$, extendible to $K$. This compression is again
zero or proportional to a pure KD-positive vector of $K$.

Fix any pure $\tau\in\hull_T$. The completely positive trace-preserving map
\[
 \mathcal R_E(X)=E^*XE+\Tr[(I-EE^*)X]\tau
\]
satisfies $\mathcal R_E(\hull_G)\subseteq\hull_T$ by the compression check,
and $\mathcal R_E(E\omega E^*)=\omega$. The inclusion
$E\hull_TE^*\subseteq\hull_G$ and trace-norm invariance under an isometry
give the ``$\leq$'' direction of \eqref{eq:isometry}. For the reverse
direction, every $\sigma\in\hull_G$ satisfies
\[
 \|\omega-\mathcal R_E(\sigma)\|_1
 =\|\mathcal R_E(E\omega E^*-\sigma)\|_1
 \leq\|E\omega E^*-\sigma\|_1,
\]
by channel contraction on Hermitian operators. Taking infima finishes
the proof.
\end{proof}

\begin{statementframe}
\begin{corollary}
Every finite abelian group with a $\Z_3^2$ subquotient admits a state in
$\E_G$ at trace-norm distance at least $111/20000$ from $\hull_G$.
In particular this holds for all noncyclic finite abelian $3$-groups.
\end{corollary}
\end{statementframe}
\begin{proof}
Compose the subgroup and quotient isometries in Theorem~\ref{thm:lift}.
A noncyclic abelian $3$-group contains a subgroup $\Z_3^2$.
\end{proof}

\begin{statementframe}
\begin{corollary}
Among finite abelian groups of order $2^a3^b$, with $a,b\geq0$,
one has $\E_G=\hull_G$ exactly for the trivial group and the cyclic groups
whose order is a power of one prime.
\end{corollary}
\end{statementframe}
\begin{proof}
For a noncyclic group, a Sylow subgroup is noncyclic and contains
$\Z_2^2$ or $\Z_3^2$. Use respectively the published example
in \cite{DLR} or Theorem~\ref{thm:example}, and then Theorem~\ref{thm:lift}.
A cyclic group with $a,b>0$ contains $\Z_6$, another published
counterexample \cite{DLR}. The remaining cyclic prime-power cases have
$\E_G=\hull_G$ by \cite{DLR}.
\end{proof}

\section{Verification materials}
Exact verification data and independent standard-library checkers accompany
this note. They reconstruct all 54 pure arrays, verify strict positivity and
the witness value, and check the diagonal-dominance certificate using exact
arithmetic. These files complement the proofs given above.

\begingroup\small
\begin{thebibliography}{99}
\setlength{\itemsep}{2pt}
\bibitem{KDreview}
D. R. M. Arvidsson-Shukur, W. F. Braasch Jr., S. De Bi\`evre, J. Dressel,
A. N. Jordan, C. Langrenez, M. Lostaglio, J. S. Lundeen and N. Yunger Halpern,
\emph{Properties and Applications of the Kirkwood--Dirac Distribution},
New J. Phys. \textbf{26}, 121201 (2024);
\href{https://doi.org/10.1088/1367-2630/ada05d}{doi:10.1088/1367-2630/ada05d};
\href{https://arxiv.org/abs/2403.18899v2}{arXiv:2403.18899v2}.

\bibitem{DLR}
S. De Bi\`evre, C. Langrenez and D. Radchenko,
\emph{The Kirkwood--Dirac representation associated to the Fourier transform
for finite abelian groups: positivity},
\href{https://arxiv.org/abs/2501.12252v2}{arXiv:2501.12252v2};
\href{https://doi.org/10.1007/s00023-025-01614-7}{doi:10.1007/s00023-025-01614-7}.

\bibitem{geometry}
C. Langrenez, D. R. M. Arvidsson-Shukur and S. De Bi\`evre,
\emph{Characterizing the geometry of the Kirkwood--Dirac-positive states},
J. Math. Phys. \textbf{65}, 072201 (2024);
\href{https://doi.org/10.1063/5.0164672}{doi:10.1063/5.0164672};
\href{https://arxiv.org/abs/2306.00086}{arXiv:2306.00086}.

\bibitem{almost}
C. Langrenez, W. Salmon, S. De Bi\`evre, J. J. Thio, C. K. Long and
D. R. M. Arvidsson-Shukur,
\emph{Almost no experiments have classical Kirkwood--Dirac representations},
Phys. Rev. A \textbf{113}, 062215 (2026);
\href{https://doi.org/10.1103/v7z4-qsz8}{doi:10.1103/v7z4-qsz8};
\href{https://arxiv.org/abs/2405.17557v2}{arXiv:2405.17557v2}.

\bibitem{convexroof}
C. Langrenez, S. De Bi\`evre and D. R. M. Arvidsson-Shukur,
\emph{Convex roofs witnessing Kirkwood--Dirac nonpositivity},
J. Phys. A: Math. Theor. \textbf{58}, 215301 (2025);
\href{https://doi.org/10.1088/1751-8121/add58f}{doi:10.1088/1751-8121/add58f};
\href{https://arxiv.org/abs/2407.04558}{arXiv:2407.04558}.

\bibitem{Spriet}
M. Spriet,
\emph{Characterizing the Kirkwood--Dirac positivity on second countable LCA groups},
\href{https://arxiv.org/abs/2507.23628v3}{arXiv:2507.23628v3}.

\bibitem{Thio}
J. J. Thio, S. Yang, N. Yunger Halpern, S. De Bi\`evre, C. H. W. Barnes
and D. R. M. Arvidsson-Shukur,
\emph{Kirkwood--Dirac Nonpositivity Is a Necessary Resource for Quantum Computing},
Phys. Rev. Lett. \textbf{137}, 080601 (2026);
\href{https://doi.org/10.1103/x819-898d}{doi:10.1103/x819-898d};
\href{https://arxiv.org/abs/2506.08092v2}{arXiv:2506.08092v2}.

\bibitem{Burkat}
J. Burkat and S. Strelchuk,
\emph{Classification of Kirkwood--Dirac positivity-preserving unitaries},
Phys. Rev. A \textbf{114}, 042446 (2026);
\href{https://doi.org/10.1103/rsy9-12fx}{doi:10.1103/rsy9-12fx};
\href{https://arxiv.org/abs/2502.11784v4}{arXiv:2502.11784v4}.

\bibitem{Cai}
L.-Y. Cai, Y.-H. Yang and Z.-J. Zheng,
\emph{Kirkwood--Dirac classical states based on discrete Fourier transform:
Representation with directed graph},
\href{https://arxiv.org/abs/2603.13863v1}{arXiv:2603.13863v1}.

\bibitem{ThioThesis}
J. Thio,
\emph{The Structure of the Kirkwood--Dirac Distribution and its Applications
in Quantum Foundations, Computation, and Information},
PhD thesis, University of Cambridge (2026), Theorem 2.2.3;
\href{https://doi.org/10.17863/CAM.135674}{doi:10.17863/CAM.135674}.
\end{thebibliography}
\endgroup
\end{document}
